% Transcription — fonds Alexandre Grothendieck, shelfmark 43, batch 4 (pages 61-80).
% Established from the facsimile archives/batches/43/batch-04.pdf (PDF page k =
% archivists' page 60+k, no cover sheet), cut from a copy of 43.pdf fetched
% through the project's own relay
% (https://grothendieck-relay.onrender.com/source/43.pdf); tiles in
% archives/tiles/43/p61-p80. Per the skill transcribe-grothendieck. The folder
% is 207 pages in eleven batches, transcribed in parallel. Batches 1-3, 5 and
% 9-11 were read for the folder's conventions: pp. 61-68 continue batch 3's
% run on adeles in several dimensions (rings A_[p,q,...] of chains of primes,
% the element Theta, pp. 44-60) and follow its notation; p. 80 ends on
% « donc » and batch 5's p. 81 opens with the Hom(U_0, G) it announces.
%
% Pass: Opus 5.5 (claude-opus-5-5), 2026-09-26 — first pass, unchecked against the pages by a human.
%
% Dating: the folder has its own inventory date, carried verbatim in \dating{}
% with no group sentence (the skill adds the group « Jacobiennes » (43-44)
% only for undated folders). No page of this batch carries a date.
%
% Copies. No page of this batch is evidently a photocopy: pp. 61-68 are ink
% on white paper, pp. 70-80 blue and blue-black ink on greyish « VELIN
% EXTRA »-watermarked paper, with the red-crayon frame of p. 70 and the
% variation of pressure and colour of originals; p. 69 is chamois card.
% All transcribed as originals.
%
% Pages skipped: none. Page 69 is a cover inscribed in his hand at the top
% right, « Jacobiennes locales »: it takes no \page marker and its words
% become the section heading (placed before p. 70, which repeats them in a
% red-crayon box), with a \note. Page 71 is almost blank: nine lines at the
% foot, written upside down relative to the archivists' number, which are a
% first redaction of a)-c) of p. 70; read turned and transcribed.
%
% Pages summarised: none. Pp. 61-63, 65-67 are sheets of scattered formulas
% and sketches (segments, lattices of chains, parallelograms); the formulas
% are given block by block and the sketches only said in \note{}s. The
% lattice of p. 61 is redrawn as a tikz-cd with its layout simplified; the
% arrays of pp. 76, 77 are laid out as on the page.
%
% His pagination: none certain. A lone « 1 » at the top left of p. 61 may
% be a leaf number (said in a \note). Numbered runs of his: a), b), c) on
% p. 70 (local cohomology H^i_A(G) for i = 0, 1, 2, « Théorie Čechiste »),
% repeated a), b), c) on p. 71 (first redaction), d) on p. 72 (H^3_A);
% sections I) « Théorie Čechiste pour la topologie de Zariski » (p. 70) and
% II « Cohomologie de Weil » (p. 76); conjectures a), b) (p. 74) and c)
% (p. 75) on J_{A/k}. No internal page reference.
%
% What the batch is about. Pp. 61-68 (end of batch 3's run on adeles for
% several dimensions): the lattice of rings A_[p], A_[p,q], A_[p,q,r] for
% p ⊃ q ⊃ r and the completions A-hat_p (p. 61); products
% A_[x_n] Theta ... Theta A_[x_0], K_c -> prod K_c', H_I(DD(A)), the
% composable-chain products A_c' x A_c -> A_c'c and their associativity
% square (p. 62); the cosimplicial-looking product prod A_[p_1...p_n],
% Theta^2 = 0, and A_[p_1..p_r] as a ring with r topologies (p. 63); the
% recursive definition A_[p] = A-hat_p, A_[p,q] = lim A-hat_p ⊗ A_q/q^n A_q
% and its residue ring (p. 64); a graded topological A-algebra A* with
% A^0 = prod O-hat_x, projectors pi_x, Theta in A^1 with Theta^2 = 0 and
% pi_x' Theta pi_x = 0 unless x' precedes x, « A libre pour ces propriétés »
% (p. 65); Hom(K_*, K_*) as a graded algebra, Hom^0 = prod O-hat_x, Theta in
% Hom^1, the graded commutator {Theta, u} (p. 66); H(x, y) = Hom(K(x), K(y))
% and the maps H_s -> H_s' for saturated chains s, an ordered set I, the
% category I-tilde and a functor K : I-tilde -> C (pp. 67-68). Pp. 70-80,
% under the cover « Jacobiennes locales »: local cohomology H^i_A(G) =
% H^i(X/X - a, G) of a regular local ring A for a commutative algebraic
% group G, computed in the Zariski (Čech) theory — G(K) in dim 0,
% G(K)/G(A) in dim 1, H^{i-1}(X - a, G) in higher dimension, reduced to the
% unipotent part via factoriality (H^1(X - a, G_m) = 0) and H^i(X - a, G_a)
% = 0 for i ≠ dim A - 1 (pp. 70-73); the local Jacobian J_{A/k} of a regular
% local ring over k, defined by Hom_{k-gr}(J_{A/k}, G) ≅ H^i_A(G), i = dim A,
% with the conjectures « Je soupçonne »: J_{K/k} an extension of Z by a
% connected group with abelian part (the Albanese), multiplicative part and
% unipotent part; J_{A/k} connected linear with multiplicative part G_m in
% dim 1; connected unipotent and projective in dim ≥ 2 (pp. 74-75); II, Weil
% cohomology: H^i(A, G) = 0 for A henselian, reduction to the residue field
% and to G/L_0, the exact sequence for H^i(A mod K, G) in dim 1, a table of
% fundamental groups pi_1(L) = pi_1(A), pi_1(K), pi_1(A mod K) (pp. 76-78);
% « Adèles et Jacobiennes généralisées des courbes » (p. 79): the spectral
% sequences Ext*(J_*, G) <= H^p(Ext^q(J_*, G)), Ext^p(H_q(J_*), G), for a
% curve X H_0(J_*) = the Rosenlicht Jacobian J^•_{X/k}, H_1(J_*) = G_m or
% 0, H*(X, G) = Ext*(J^•_{X/k}, G) for X affine, the long exact sequence
% for X complete, and, boxed, Ext*(J^•_{K/k}, G) ≅ H*(K, G) for K = k(X);
% p. 80, computation of J^•_{A/k} for a DVR with residue field of
% transcendence degree 1: coprod S*_L = coprod Z_L x coprod G_m,L x coprod
% S'_L, the connected unipotent U = coprod S'_L, and Hom_{k-gr}(U, G)
% identified, through Rosenlicht, with restricted Hom(J^•_{L/k} x E^1, G).
%
% Notation fixed once for the batch: \mathfrak{p}, \mathfrak{q}, \mathfrak{r},
% \mathfrak{m} for his gothic letters, \mathcal{A} for his cursive A, \Theta,
% \pi_x, \coprod for his upside-down product, as in batch 3; J plain for his
% J (J_{A/k}, J_{K/k}, J_*), J^\bullet where he dots it (the Rosenlicht
% Jacobian, pp. 79-80); \mathcal{H} for his cursive H (cohomology of
% arbitrary principal bundles, p. 72) against the straight H of the Zariski
% theory; \underline{O}_x, \underline{R}_x, \underline{D}_x (p. 72);
% \mathbb{G}_m, \mathbb{G}_a; \mathbb{S} for his barred S (p. 80, as batch 2);
% \mathcal{U} for his looped U (batch 5 reads the same letter as U_0 on
% p. 81); L_0 for the largest connected linear subgroup; A mod K for the
% relative cohomology of pp. 77 (reading of « mod » doubtful).
%
% Readings to check first: the whole prose of pp. 63 (foot), 64 and 70-72,
% which is fast drafting — the formulas carry, the connecting words often
% only as \uncertain{}; the lattice and the labels of p. 61; the marginal
% notes of pp. 70, 74 (left, oblique, heavily overwritten) and 79 (right);
% the exponents of H_A on pp. 71 (c) and 72 (d); « biprécède » on p. 67;
% the sign of Theta u - u Theta on p. 66 (transcribed as read, flagged);
% the word read « mod » on p. 77.
% Apparatus: about 197 \ill{} and 871 \uncertain{} in the body.

\documentclass[11pt,a4paper]{article}
% Le préambule commun à toutes les transcriptions du fonds Grothendieck.
%
% Il tient en une idée : **l'incertitude se compose, elle ne se résout pas.**
% Une transcription qui lisse un mot douteux a détruit la seule chose pour
% laquelle on la faisait — un lecteur doit pouvoir savoir, à chaque endroit,
% ce qui était sur le feuillet et ce qui a été ajouté. Les six macros qui
% suivent sont donc l'appareil critique, et non de la décoration.
%
% Les mêmes commandes sont reconnues par `scripts/render.mjs`, qui produit la
% vue de lecture du site. Une macro ajoutée ici doit l'être aussi là-bas,
% faute de quoi le rendu échouera bruyamment — ce qui est voulu : un rendu
% silencieusement faux est pire qu'un rendu absent.

\ProvidesPackage{grothendieck}[2026/08/08 Transcriptions du fonds Grothendieck]

\usepackage{amsmath,amssymb,amsthm}
\usepackage{mathrsfs}
\usepackage{stmaryrd}
% Les diagrammes commutatifs sont partout dans ce fonds : c'est la langue
% même de la géométrie algébrique de Grothendieck, et une transcription qui
% les rendrait en prose ne transcrirait rien.
\usepackage{tikz-cd}
% Français par défaut — c'est la langue des feuillets — mais l'anglais est
% chargé aussi : la traduction et le résumé sont des documents anglais, et la
% typographie française (espaces avant les deux-points) y serait une faute.
% Ces documents basculent par \selectlanguage{english} en tête de corps.
\usepackage[english,french]{babel}
\usepackage{fontspec}
\usepackage{geometry}
\geometry{a4paper,margin=25mm,marginparwidth=28mm}
\usepackage{marginnote}
\usepackage{xcolor}
% ulem, not soul. `soul`'s \st and \ul are built on a character-by-character
% scan that XeLaTeX's font machinery breaks: the document fails with an
% undefined control sequence rather than degrading. ulem does the same two jobs
% and is Unicode-engine safe. `normalem` stops it hijacking \emph, which the
% transcriptions use for Grothendieck's own emphasis.
\usepackage[normalem]{ulem}
\usepackage{hyperref}
\hypersetup{colorlinks=true,linkcolor=black,urlcolor=black,pdfborder={0 0 0}}

% --- Métadonnées du lot -----------------------------------------------------
% Reprises telles quelles par le script de rendu, qui les affiche en tête de
% la vue de lecture. Elles fixent ce que le fichier prétend couvrir : sans
% elles, une transcription flotte sans point d'attache dans un fonds de seize
% mille feuillets.

\newcommand{\folder}[1]{\gdef\grothFolder{#1}}
\newcommand{\batch}[1]{\gdef\grothBatch{#1}}
\newcommand{\leaves}[2]{\gdef\grothFirst{#1}\gdef\grothLast{#2}}
\newcommand{\foldertitle}[1]{\gdef\grothFoldertitle{#1}}
\gdef\grothFolder{?}\gdef\grothBatch{?}\gdef\grothFirst{?}\gdef\grothLast{?}\gdef\grothFoldertitle{}

% --- Le résumé --------------------------------------------------------------
%
% Il ouvre la lecture modernisée, sous le titre « Résumé », et il est composé
% en petit. Ce n'est pas un ornement : c'est le seul passage du document qui
% ne soit pas des mathématiques, et un lecteur doit pouvoir le reconnaître
% avant de l'avoir lu — puis le sauter s'il connaît déjà le sujet, ou s'y
% arrêter s'il ne le connaît pas. Le filet qui le suit marque la reprise du
% corps.

\newenvironment{resume}
  {\small\setlength{\parskip}{0.45em}}
  {\par\medskip\hrule\medskip}

% --- Mots-clés ---------------------------------------------------------------
%
% En anglais, en clôture du résumé de la lecture modernisée : le vocabulaire
% moderne sous lequel chercher ce que les feuillets construisent. Le manifeste
% du site (`npm run manifest`) les extrait pour étiqueter la cote entière —
% cette ligne est la source unique des tags, il n'y a pas de fichier à côté.
\newcommand{\keywords}[1]{\par\smallskip\noindent
  {\footnotesize\textsc{Keywords} --- \textit{#1}}\par}

% --- L'appareil critique ----------------------------------------------------

% Le numéro de feuillet, en marge. C'est l'ancre qui permet de régler un
% désaccord sur une formule en regardant *un* feuillet plutôt qu'une chemise
% de six cents. Le site s'en sert aussi pour tourner les pages du fac-similé
% quand on fait défiler la transcription.
\newcommand{\leaf}[1]{%
  \marginnote{\footnotesize\color{gray!70}\textsf{#1}}%
  \label{leaf:#1}\ignorespaces}

% Illisible : on ne devine pas. Un mot inventé qui se lit comme les autres est
% le pire résultat possible.
\newcommand{\ill}{\textcolor{red!60!black}{[\,\ldots\,]}}

% Lecture douteuse : le mot est donné, et signalé comme incertain.
\newcommand{\uncertain}[1]{\uline{#1}}

% Ajout de l'éditeur — une lettre manquante, un mot sous-entendu.
\newcommand{\add}[1]{\textcolor{blue!50!black}{[#1]}}

% Biffé par Grothendieck lui-même. Ce qu'il a rayé fait partie du document :
% une rature dit souvent où la pensée a changé de direction.
\newcommand{\struck}[1]{\sout{#1}}

% Note du transcripteur, sur l'état matériel du feuillet.
\newcommand{\note}[1]{\par\smallskip{\small\color{gray!60!black}\textit{[#1]}}\par\smallskip}

% Note marginale de Grothendieck, distincte du corps du texte.
\newcommand{\marginal}[1]{\par\smallskip\noindent\hspace{1em}%
  {\small\color{gray!60!black}#1}\par\smallskip}

% --- Titre ------------------------------------------------------------------

\AtBeginDocument{%
  \begin{center}
    {\large\bfseries Fonds Alexandre Grothendieck --- cote n\textsuperscript{o}~\grothFolder}\\[2pt]
    {\itshape\grothFoldertitle}\\[4pt]
    {\small feuillets \grothFirst--\grothLast\ (lot \grothBatch)}\\[2pt]
    {\footnotesize\color{gray!60!black}Transcription établie d'après le fac-similé numérisé
     par l'Université de Montpellier.\\
     Les lectures incertaines sont soulignées, les illisibles notées [\,\ldots\,],
     les ajouts entre crochets.}
  \end{center}
  \vspace{1em}}

\endinput


\folder{43}
\batch{4}
\pages{61}{80}
\dating{[à partir de 1961-vers 1968]}
\watermark{Édition de démonstration}
\foldertitle{Dualité des groupes. Jacobiennes généralisées etc : notes et copies de notes manuscrites (s.d.)}

\begin{document}

\page{61}
\note{feuille de schémas épars, sans prose suivie ; elle semble prolonger
les pages 51--56 (anneaux $A_{[\mathfrak{p}, \mathfrak{q}]}$ attachés aux
chaînes d'idéaux premiers) et 59--60 (l'élément $\Theta$). En haut à gauche,
un « 1 » isolé, peut-être un numéro de feuillet. Les schémas sont donnés
dans l'ordre de la page}

\begin{tikzcd}
\hat{A}_{\mathfrak{p}} \arrow[r, "{\mathfrak{q}\hat{A}_{\mathfrak{p}}}"] & A_{[\mathfrak{p}, \mathfrak{q}]} \\
A \arrow[u] \arrow[r] & \hat{A}_{\mathfrak{q}} \arrow[u]
\end{tikzcd}

\note{au-dessus de la flèche horizontale, « l'\ill{} \ill{} \ill{} » ; au
centre du carré « \uncertain{commutatif} » ; à droite de la flèche
verticale, « \ill{} \uncertain{local} »}

\begin{tikzcd}[column sep=small]
 & A_{[\mathfrak{p}, \mathfrak{q}, \mathfrak{r}]} & & \\
A_{[\mathfrak{p}, \mathfrak{q}]} \arrow[ur] & & A_{[\mathfrak{q}, \mathfrak{r}]} \arrow[ul] & \\
A_{[\mathfrak{p}]} \arrow[u] & A_{[\mathfrak{q}]} \arrow[ul] \arrow[ur] & & A_{[\mathfrak{r}]} \arrow[ul] \\
 & A_{\mathfrak{p}} \arrow[ul] \arrow[u] \arrow[urr] & &
\end{tikzcd}

\note{le treillis est redessiné en simplifiant sa disposition ; la flèche
de $A_{\mathfrak{p}}$ vers $A_{[\mathfrak{r}]}$ est légendée, dans un ovale,
« \uncertain{complétion} \uncertain{par} $\mathfrak{r}$ » ; dessous :
$\mathfrak{p} \supset \mathfrak{q} \supset \mathfrak{r}$}

À droite :
\[
\pi_{x_n} \Theta \cdots \Theta \pi_{x_2} \Theta \pi_{x_1}
\qquad
\mathfrak{p} \supset \mathfrak{q} \supset \mathfrak{r}
\]
\note{sous le produit, une flèche horizontale épaisse, de gauche à droite}
\[
\begin{array}{ll}
A_{[\mathfrak{p}, \mathfrak{q}]} & \\
\uparrow \text{\ill{} \uncertain{plat} \ill{} \ill{}} & \\
A_{\mathfrak{q}} & A_{\mathfrak{q}} / \mathfrak{q} A_{\mathfrak{q}}
\end{array}
\]
\marginal{il y a \ill{} \ill{} \ill{} de dim.\ 1, \ill{} $A/\mathfrak{q}$,
\ill{} \ill{} \ill{} \ill{} \ill{} \ill{}}
\[
A_{[\mathfrak{p}, \mathfrak{q}]} \otimes_{A_{\mathfrak{q}}} A_{\mathfrak{r}} / \mathfrak{r}^{n} A_{\mathfrak{r}}
\qquad
A_{[\mathfrak{p}, \mathfrak{q}]} / \mathfrak{r}^{n} A_{[\mathfrak{p}, \mathfrak{q}]} \otimes A_{\mathfrak{r}}
\]
\note{sous le second terme, souligné : « dim.\ 1 », puis
« \uncertain{complétion} \uncertain{de} $A/$ » ; au centre de la page, un
petit carré vide marqué d'un point}

Au bas de la page :

\begin{tikzcd}[column sep=small]
A_{[\mathfrak{p}_1 \cdots \mathfrak{p}_{n-1}]} \arrow[r] & A_{[\mathfrak{p}_1, \ldots, \mathfrak{p}_n]} \\
\hat{A}_{\mathfrak{p}_1} \cdots \hat{A}_{\mathfrak{p}_i} \cdots \arrow[ur] & \hat{A}_{\mathfrak{p}_n} \arrow[u]
\end{tikzcd}

\note{la flèche du haut porte une lettre biffée ; au-dessus des
$\hat{A}_{\mathfrak{p}_i}$, « $A_{[\mathfrak{p}_1}$ » biffé ; à droite,
« $A_{[\mathfrak{p}_1, \ldots,}$ » inachevé. En bas, trois croquis sans
légende lisible : un réseau de courbes autour d'un point ; une flèche
verticale de $\mathcal{C}$ vers « \uncertain{Ens.} \ill{} » ; et
$x_0 \longrightarrow x_n$ surmonté d'un arc, légendé « \uncertain{chaîne}
\uncertain{saturée} »}

\page{62}
\note{feuille de formules éparses, dans la suite de la précédente et des
pages 59--60 (anneaux $A_c$ des chaînes saturées, élément $\Theta$) ; les
blocs sont donnés de haut en bas}
\[
A_{[x_n]} \Theta A_{[x_{n-1}]} \cdots \Theta A_{[x_1]} \Theta A_{[x_0]}
\]
\note{sous la formule, une flèche vers la gauche ; à droite, trois petits
segments gradués (chaînes dessinées), un signe \ill{} et}
\[
K_c \longrightarrow \prod_{c'} K_{c'}
\qquad
\mathrm{Hom}(K_{*}, K_{*})
\qquad
H_{I}\bigl(DD(A)\bigr) = \underline{H}^{**}(A)
\qquad
A_x
\]
\note{un trait horizontal traverse ensuite la page ; sur ce trait, un
petit croquis de courbe marquée d'un point}
\[
(u_1 - u_2)^{2} = u_1 u_1 - u_1 u_2 - u_2 u_1
\qquad
\sum 1_{x_0, x_1}
\qquad
u_1
\]
\note{le développement est inachevé à droite du trait}
\[
A_{c'} \times A_{c} \longrightarrow A_{c'c}
\quad \text{si \emph{$c'$ et $c$ composables} (\uncertain{bout} à \uncertain{bout})}
\]
\[
c = (x_0, \ldots, x_n), \qquad c' = (x_n, \ill{}, x_m)
\]
\[
\text{\ill{}} \qquad
\begin{array}{l}
A_{c'} \longrightarrow A_{c'c} \\
A_{c} \longrightarrow A_{c'c}
\end{array}
\]
\note{suit une ligne soulignée : « \struck{\ill{}} \emph{\ill{} \ill{}
\uncertain{produit} \ill{}} »}

\begin{tikzcd}
A_{c''} \times A_{c'} \times A_{c} \arrow[r] \arrow[d] & A_{c''c'} \times A_{c} \arrow[d] \\
A_{c''} \times A_{c'c} \arrow[r] & A_{c''c'c}
\end{tikzcd}

\note{dans le coin inférieur gauche, un signe biffé entre $A_{c''}$ et
$A_{c'c}$, lu $\times$}
\[
f, g, h \longmapsto \bigl(i_{c'', c''c'}(f)\, i_{c', c''c'}(g),\ h\bigr)
\longmapsto
i_{c''c', c''c'c}(\,\cdots\,)\, i_{c, c''c'c}(h)
= i_{c'', c''c'c}(f)\, i_{c', c''c'c}(g)\, i_{c, c''c'c}(h)
\]
\note{la formule est disposée en deux étages reliés par une flèche
verticale et une accolade ; un signe biffé la précède. Au-dessous, à gauche,
trois segments consécutifs légendés $c$, $c'$, $c''$, et « $A_c$ »}

\page{63}
\note{feuille de formules et de croquis ; un trait horizontal la partage
en deux. Les blocs sont donnés de haut en bas}
\[
\underbrace{\prod_{\mathfrak{p}_1 \supset \mathfrak{p}_2 \supset \cdots \supset \mathfrak{p}_n} A_{[\mathfrak{p}_1, \mathfrak{p}_2, \ldots, \mathfrak{p}_n]}}_{\mathcal{A}^{n}}
\longleftarrow
\prod_{\mathfrak{q}_1 \supset \cdots \supset \mathfrak{q}_{n-1}} A_{[\mathfrak{q}_1, \ldots, \mathfrak{q}_{n-1}]}
\]
\marginal{\uncertain{cosimplicial} ?!}
\[
\mathfrak{p}_1 \supset \mathfrak{p}_2 \supset \cdots \supset \mathfrak{p}_n
\qquad
x_1 \leq x_2 \leq \cdots \leq x_n
\qquad
A_{\mathfrak{p}, \mathfrak{q}}
\qquad
\prod_{\mathfrak{q} \subset \mathfrak{p}} A_{[\mathfrak{q}, \mathfrak{p}]}
\qquad
\Theta^{2} = 0
\]
\note{sous les chaînes, deux croquis de segments gradués parcourus par
des flèches vers la gauche, un petit parallélogramme marqué d'un point, et
un « $A$ » isolé}
\[
\Sigma\, 1_{\mathfrak{q}, \mathfrak{p}}
\qquad
[\mathfrak{p}_1 \supset \cdots \supset \mathfrak{p}_n]
\qquad
\coprod_{\mathfrak{q}} 1_{\mathfrak{q}, \mathfrak{p}_1}
\qquad
\downarrow 1_{\mathfrak{q}\, \mathfrak{p}}
\]
\[
\prod_{\mathfrak{q}} [\mathfrak{q} \supset \mathfrak{p}_1 \supset \cdots \supset \mathfrak{p}_n]
\qquad\qquad
\prod_{\mathfrak{q}} [\mathfrak{p}_1 \supset \cdots \supset \mathfrak{p}_n \supset \mathfrak{q}]
\]
\note{des lignes courbes et des flèches relient ces termes : la
multiplication par $\Theta = \Sigma\, 1_{\mathfrak{q}, \mathfrak{p}}$ à
gauche allonge la chaîne par le haut, à droite par le bas ; le dessin est
résumé ici}


$A_{[\mathfrak{p}_1, \ldots, \mathfrak{p}_r]}$ \uncertain{anneau}
\uncertain{local} \uncertain{muni} de $r$ \uncertain{topologies}
(\uncertain{non} \uncertain{adiques} en \uncertain{général})
\uncertain{différentes}, $T_1, \ldots, T_r$ (\uncertain{pas}
\uncertain{nécessairement} \uncertain{comparables})
\uncertain{métrisables} \struck{\ill{}} et \uncertain{complètes}.
\note{à droite : $\hat{A}_{\mathfrak{p}} / \mathfrak{q}^{n} \hat{A}_{\mathfrak{q}}$
et $\hat{A}_{\mathfrak{q}}$ ; à gauche, un croquis de quadrilatère marqué
d'un point}
\ill{} \uncertain{une} \uncertain{suite} \uncertain{filtrante}
\uncertain{tend} \uncertain{vers} $0$ \uncertain{pour} $T_i$, \ill{}
\ill{} \ill{} \ill{} \ill{} \ill{} \uncertain{voisinage} \uncertain{de}
$0$ $V_i$, \uncertain{pour} $T_{i-1}$,
\note{la page s'arrête sur cette virgule}

\page{64}
\note{feuille de formules ; en haut, un petit croquis de deux crochets
chevauchants au-dessus de « $\circ$ \quad $\mathfrak{r}$ \quad
$\mathfrak{m}$ »}
\[
\hat{A}_{\mathfrak{r}} \otimes_{A} K
\qquad
\widehat{\hat{A} \otimes A_{\mathfrak{r}}}
\qquad
\widehat{A/\mathfrak{r}^{n}A} \otimes A
\qquad
\underline{A}_{\mathfrak{r}}
\qquad
A/\mathfrak{r}^{n}A
\]
\note{des courbes relient $\underline{A}_{\mathfrak{r}}$ aux deux derniers
termes}
\[
\begin{array}{l}
A/\mathfrak{m} = B_0 = L_0 \\
\widehat{A/\mathfrak{r}} = B_1 \subset L_1 = \text{\uncertain{anneau} \uncertain{local} \uncertain{des} \uncertain{fractions} \uncertain{de} } \widehat{A/\mathfrak{r}} = B_1 \\
\varprojlim L_1 \qquad = \hat{A}_{\mathfrak{r}} \otimes_{A} A_{\mathfrak{r}} / \mathfrak{r} A_{\mathfrak{r}}
\end{array}
\]
\note{dans la première ligne, un signe biffé avant $B_0$. À gauche, un croquis de deux courbes se coupant, l'une marquée d'un
point}
\[
A \qquad \mathfrak{p}_1 \supset \mathfrak{p}_2 \supset \cdots \supset \mathfrak{p}_n
\qquad
A_{[\mathfrak{p}_1, \mathfrak{p}_2, \ldots, \mathfrak{p}_n]}
\]
\note{au-dessus de $\mathfrak{p}_1 \supset \mathfrak{p}_2$ et de
$\supset \mathfrak{p}_n$, deux petites marques en chapeau}
\uncertain{définition} \uncertain{par} \uncertain{récurrence}
\uncertain{sur} $n \geq 1$.
\[
A \to \hat{A}_{\mathfrak{p}} \quad \text{\ill{} \uncertain{local} \uncertain{complété}}
\qquad
A_{[\mathfrak{p}]} = \hat{A}_{\mathfrak{p}}
\]
\uncertain{si} $n = 1$ ; c'\uncertain{est} \uncertain{un}
\struck{\ill{}} \uncertain{anneau} \struck{\ill{}} $A$ \uncertain{local}
\uncertain{noeth.}\ \uncertain{et} \uncertain{complet} \uncertain{et}
\uncertain{une} \uncertain{algèbre} \uncertain{sur} $A$.
\[
A_{[\mathfrak{p}_1 \cdots \mathfrak{p}_n]}
\]
\uncertain{déjà} \uncertain{défini}, c'\uncertain{est} \ill{} \ill{}
\ill{} \uncertain{local} \uncertain{noethérien} \uncertain{séparé} et
\uncertain{complet}, et \uncertain{une} \uncertain{algèbre} \uncertain{sur}
$A$.
\[
A_{[\mathfrak{p}, \mathfrak{q}]} = \varprojlim_{n} \hat{A}_{\mathfrak{p}} \otimes_{A} A_{\mathfrak{q}} / \mathfrak{q}^{n} A_{\mathfrak{q}}
\qquad (\mathfrak{p} \supset \mathfrak{q})
\qquad
\widehat{A_{\mathfrak{p}}/\mathfrak{q}A_{\mathfrak{p}}} \otimes
\]
\ill{} \ill{} \uncertain{complet} \uncertain{dont} \uncertain{le}
\struck{\ill{}} \uncertain{corps} \uncertain{résiduel} $K(\mathfrak{p}, \mathfrak{q})$
\note{dans la formule précédente, un signe biffé avant $\varprojlim$}
\uncertain{est} \uncertain{l'anneau} \uncertain{total} \uncertain{des}
\uncertain{fractions} \uncertain{de}
\[
\widehat{A_{\mathfrak{p}}/\mathfrak{q}A_{\mathfrak{p}}} = \hat{A}_{\mathfrak{p}} / \mathfrak{q} \hat{A}_{\mathfrak{p}}
\]
(\uncertain{complété} d'un \uncertain{anneau} \uncertain{local}
\uncertain{intègre} \uncertain{noeth.}\ \uncertain{de} \uncertain{dim.}\ 1)

\page{65}
\note{feuille de formules et de croquis, puis, au bas, une définition
rédigée ; les blocs épars du haut sont donnés dans l'ordre de la page. En
tête, à gauche, un $\mathcal{A}$ cursif isolé}
\struck{$\Theta^{2} = 0$}
\[
A \longrightarrow \mathcal{B}
\qquad
\Theta_{x x'} =
\qquad
\Theta
\qquad
\Theta^{2} = 0
\]
\note{après le second $\Theta$, un signe biffé ; sous $\Theta^2 = 0$,
« $A \otimes B$ » \uncertain{biffé}}
\[
\pi_{x'} \Theta \pi_{x} = \Theta_{x x'}
\qquad
x_1 \geq x_2 \geq x_3 \geq \cdots \geq x_n \geq x_{n+1} \geq x_{n+2}
\qquad
\Theta_{x_1, x_2}\ \Theta
\]
\note{entre les deux, une ligne biffée : \struck{$x_1 \subset x_2 \subset X$}}
\note{à gauche, un croquis de parallélogramme ; à droite, deux segments
gradués, dont l'un a un intervalle épaissi}
\[
\pi_{x_m} \Theta \cdots \pi_{x_{n+2}} \Theta \pi_{x_{n+1}}\
\underbrace{\varphi_{x_n x_{n+1}}}\
\underbrace{\pi_{x_n} \Theta \cdots \pi_{x_3} \Theta \pi_{x_2} \Theta \pi_{x_1}}
\]
\[
\varphi_{x} \quad \varphi_{x'}
\qquad
\left\lbrace
\begin{array}{l}
\Theta\ \text{\uncertain{$A$-linéaire}} \\
\Theta^{2} = 0 \\
\Theta\ \text{\uncertain{de} \uncertain{degré} 1}
\end{array}
\right.
\qquad
\Theta
\]
\note{autour, trois croquis : un parallélogramme aux sommets marqués
$\mathfrak{m}$ et $\mathfrak{p}$, un carré en pointillé parcouru de flèches,
et un segment gradué $\circ$, $\mathfrak{r}$, $\mathfrak{m}$ avec une flèche
oblique ; sous le premier, le mot « Hom »}


$\mathcal{A}$ \struck{anneau gradué topologique,}
\note{$\mathcal{A}$ surmonte cette ligne biffée ; la ligne suivante la
remplace}

$\mathcal{A}^{*}$ \struck{\ill{}} \uncertain{algèbre}
\uncertain{topologique} (\uncertain{séparée}, \uncertain{complète})
\uncertain{graduée} \uncertain{sur} $A$, \uncertain{à} \uncertain{degrés}
$\geq 0$, avec
\note{« (séparée, complète) » est ajouté au-dessus de la ligne ; au lieu
de « $\geq 0$ », un premier signe est surchargé}
\[
\left\lbrace
\begin{array}{l}
\mathcal{A}^{0} = \prod_{x \in X} \hat{\mathcal{O}}_{x}
\quad \text{\uncertain{d'où} \uncertain{des}}\ \pi_{x} \in \mathcal{A}_0 \quad (x \in X) \\
\qquad \text{(\uncertain{comme} \uncertain{A-algèbre} \uncertain{topologique})} \\[4pt]
\Theta \in \mathcal{A}^{1}, \quad \text{avec}\ \underline{\Theta^{2} = 0}, \quad
\underline{\pi_{x'} \Theta \pi_{x} = 0\ \text{si}\ x'\ \text{\uncertain{ne} \uncertain{précède} \uncertain{pas}}\ x}
\end{array}
\right.
\]
\note{un signe biffé suit « $\mathcal{A}^{0} =$ » ; dans $(x \in X)$, le $X$ est écrit sur un $A$ ; un mot biffé précède « si »}

$\mathcal{A}$ \uncertain{libre} pour ces propriétés \ldots
\note{la page s'arrête sur ces points}

\page{66}
\note{feuille de formules ; en tête, deux lignes biffées :
\struck{Hom} et \struck{$H^{*}(A) = \varinjlim$}}
\[
\underbrace{\mathrm{Hom}(K_{*}, K_{*})}_{\substack{\text{\uncertain{$A$-algèbre} \uncertain{graduée}} \\ \text{\uncertain{par} \uncertain{la} \uncertain{composition} \uncertain{des} \uncertain{Hom}}}}
= \prod_{x} \mathrm{Hom}\bigl(K(x), K\bigr)
\supset \prod_{x} \coprod_{y} \mathrm{Hom}\bigl(K(x), K(y)\bigr)
\]
\note{dans le premier $\mathrm{Hom}$, un indice biffé après le premier
$K$ ; à droite : $K(x) \to \coprod_{x'} K(x')$}
\[
\prod_{x} \mathrm{Hom}\bigl(K(x), K(x)\bigr)
\qquad
\prod_{x}
\]
\note{ligne biffée : \struck{$\mathrm{Hom}^{0}(K$}}
\[
\mathrm{Hom}^{0}(K, K)^{0} = \prod_{x} \hat{\mathcal{O}}_{x}
\qquad
\text{\ill{} \uncertain{composantes}} \quad (\pi_{x})
\]
\note{$\pi_x$ est entouré}
\[
\underset{\text{(\uncertain{l'application} \uncertain{nulle})}}{\Theta}
\in \mathrm{Hom}(K, K)^{1} = \prod_{x} \mathrm{Hom}\Bigl(K(x), \coprod_{\underline{\underline{x'\ \text{précède}\ x}}} K_{x'}\Bigr)
\simeq \prod_{\substack{\text{\uncertain{restreint}} \\ x'\ \text{\uncertain{précède}}\ x}} \hat{\mathcal{O}}_{x, x'}
\]
\[
\pi_{x'} \Theta \pi_{x} \qquad \hat{\mathcal{O}}_{x, x'}
\]
\note{avant $\pi_{x'}$, un signe biffé ; à droite, un groupe biffé et
surchargé, où se lisent $\Theta$ et $\mathcal{O}_{x'}$}
\[
\begin{array}{cc}
\mathcal{O}_{x'} & \longrightarrow \ \mathcal{O}_{x} \\
\| & \| \\
A & A_{\mathfrak{p}}
\end{array}
\qquad\qquad
\Sigma\, \varphi_{x'} \Theta \varphi_{x}
\qquad \pi_{x'}
\]
\note{au centre, un croquis : deux cercles sécants et un point}
\[
A_{\mathfrak{p}} \rightsquigarrow I_{A}
\qquad
\hat{A} \otimes_{A} A_{\mathfrak{p}} / \mathfrak{p}^{n} A_{\mathfrak{p}}
\qquad
D(A_{\mathfrak{p}}/\mathfrak{p}^{n}A_{\mathfrak{p}}) \to I_{A_{\mathfrak{p}}/\mathfrak{p}^{n}A_{\mathfrak{p}}}
\qquad
\hat{A} \otimes_{A} \hat{A}_{\mathfrak{p}}
\qquad
\widehat{A_{\mathfrak{p}}}
\qquad
A_{\mathfrak{p}, \mathfrak{p}'}
\]
\note{dans le deuxième terme, le signe $\otimes$ est surchargé}
\[
\Gamma(u\, v) =
\qquad
[\Theta, x]
\qquad
\mathrm{Hom}(K, K)
\]
\[
\lbrace \Theta, u \rbrace = \Theta u + (-1)^{\deg \Theta \deg u} u \Theta = \Theta u + (-1)^{\deg u} u \Theta
\]
[\uncertain{si} $\deg u = 0$, $\deg \Theta = 1$, \ill{} \ill{}
$\Theta u - u \Theta$]
\note{le signe entre crochets semble contredire la formule qui précède
(on attendrait $\Theta u + u \Theta$ pour $\deg u = 0$) ; il est transcrit
tel qu'on le lit}

\page{67}
\note{feuille de formules et de croquis, dans la suite de la
précédente ; les blocs sont donnés de haut en bas}
\[
K(x) \longrightarrow \coprod_{y} K(y)
\qquad
K(x) \longrightarrow \prod_{x'\ \text{\uncertain{précède}}\ x}
\]
\[
\left\lbrace
\begin{array}{c}
x \\ \downarrow \\ x' \\ \mid \\ x''
\end{array}
\right.
\qquad
\begin{array}{c}
K(x) \longrightarrow \coprod_{x'\ \text{précède}\ x} K(x') \\
\downarrow \\
\coprod_{x''\ \text{biprécède}\ x} K(x'')
\end{array}
\]
\[
\underbrace{K(x) \longrightarrow \coprod_{\substack{x' \\ \text{\uncertain{pts} \uncertain{intermédiaires}} \\ \text{\uncertain{entre} } x \text{ et } x''}} K(x') \longrightarrow K(x'')}_{\text{\uncertain{nul}}}
\]
\note{ligne biffée : \struck{$\mathrm{Hom}\bigl(K(x), K(y)\bigr)$}}
\[
\mathrm{Hom}(K, K) = \prod_{x} \mathrm{Hom}\bigl(K(x), K\bigr)
= \prod_{x} \prod_{y}^{\text{\uncertain{restreint}}} \underbrace{\mathrm{Hom}\bigl(K(x), K(y)\bigr)}_{H(x, y)}
\]
\[
H(x, x)\ \text{\uncertain{un} \uncertain{anneau}}
\qquad
H(x
\qquad
K \rightsquigarrow I
\qquad
\mathrm{Hom}(K,
\]
\[
K(\mathfrak{p}) \longrightarrow I
\qquad
\mathrm{Hom}\bigl(K, D(A)\bigr)\ \text{(\uncertain{anneau} \uncertain{des} \uncertain{fractions})}
\qquad
\mathfrak{q} \supset \mathfrak{p}
\]
\note{le bas de la page est occupé par des croquis : un segment épaissi
relié par des flèches à $K(x)$, $K(x')$, $K(y)$, $K(y')$ ; un segment
gradué parcouru d'arcs ; deux graphes de points reliés par des flèches
étiquetées $d$, $d'$, $d_s'$ et $\Theta$ (entouré), le second portant les
poids 1, 2, 3 ; de petits segments légendés « $d$ \uncertain{gauche} »,
« \uncertain{gauche} $d$ » ; un segment gradué sous
$\Theta\, \Theta\, \Theta \quad \Theta\, \Theta$ ; et un dernier
croquis légendé « \uncertain{gauche} $\Theta$ »}

\page{68}

\begin{tikzcd}[column sep=small]
 & K(x) \arrow[r] & K(y) \arrow[dr] & \\
K(x') \arrow[ur] \arrow[rrr] & & & K(y')
\end{tikzcd}

\note{au-dessus de la flèche du haut : $x \geq y$ ; sous la flèche du
bas : $x' \geq x \geq y \geq y'$, le $\geq$ central sous une accolade}
\[
H(x, y) \longrightarrow H(x', y')
\]
\uncertain{chaque} \uncertain{fois} \uncertain{qu'on} \uncertain{a}
\uncertain{donné} \uncertain{une} \uncertain{chaîne} \uncertain{saturée}
\uncertain{joignant} $x$ \uncertain{à} $x'$ et \uncertain{une}
\uncertain{autre} \uncertain{joignant} $y'$ \uncertain{à} $y$.
\[
y' \rightsquigarrow y \longrightarrow x \rightsquigarrow x'
\]
\uncertain{Si} : \struck{\ill{}} \uncertain{chaîne} \uncertain{saturée}
$s$ \uncertain{joignant} $x$ \uncertain{à} $y$, \uncertain{on}
\uncertain{associe} $H_s = H(x, y)$, \uncertain{alors} si
$s \leq s'$, \uncertain{alors} $H_s \to H_{s'}$ ;
\struck{d'ailleurs $H_{s'}$ est \uncertain{engendré} par les images des
$H_s$, pour $s' \leq s$, s'\uncertain{identifie} à $s$}
\note{le passage biffé est d'une lecture incertaine ; un « $i$ » est
écrit au-dessus de « $H_{s'}$ »}
\[
H_s \quad s\ \text{de longueur}\ 0
\qquad
H_s \quad s\ \text{de longueur}\ 1
\]
\[
\Bigl(\coprod_{s \in C_2(x, y)} H_s\Bigr) \Bigl(\coprod_{t \in C_{1}(x, y)} H_t\Bigr)
\qquad
\uparrow
\qquad
\underset{s \in C_0(x, y)}{H}
\]
\note{la flèche verticale monte de $H$, indexé par $s \in C_0(x, y)$,
vers le produit ; l'indice du second coproduit est d'une lecture
incertaine. À gauche, des croquis de segments et de deux chaînes joignant
deux points, entourées d'accolades}
\[
K_n = \coprod H_s
\qquad
\prod_{x} H\bigl(x, \textstyle\coprod y\bigr)
\qquad
\prod_{x} \prod_{y\ \text{\uncertain{local}}} H(x, y)
\]
\note{à droite, une suite de flèches doubles et simples,
$\rightrightarrows \rightrightarrows \to \to$}

$I$ \uncertain{ensemble} \uncertain{ordonné} « \uncertain{discret} »
\uncertain{en} \ill{} \ill{}

$\tilde{I}$ \uncertain{catégorie} \uncertain{des} \uncertain{flèches}
\uncertain{composées}

$K$ \struck{\ill{}} \uncertain{foncteur} \uncertain{de} $\tilde{I}$ dans $\mathcal{C}$

\section{Jacobiennes locales}
\note{titre pris sur la page 69, une couverture de carton chamois
inscrite de sa main, en haut à droite, « Jacobiennes locales » ; la
page 70 répète ces mots, encadrés au crayon rouge, en tête de page}

\page{70}
\uncertain{Tout} \uncertain{dépend} \uncertain{des} \uncertain{quelles}
\uncertain{théories} \ill{} \ill{} \uncertain{place} :
\marginal{\emph{Jacobiennes locales}, encadré au crayon rouge}

I) \emph{Théorie Čechiste} \uncertain{pour} \uncertain{la}
\uncertain{topologie} \uncertain{de} Zariski. (\uncertain{celle}
\uncertain{où} le $H^1$ \uncertain{classifie} les \struck{\ill{}}
\uncertain{fibrés} \emph{principaux} \emph{\uncertain{loc.}\
\uncertain{triviaux}}). \ill{} \ill{} \ill{} \ill{} \ill{}
\uncertain{exacte} \ill{} \ill{}
$0 \to G' \to G \to G'' \to 0$ \ill{} \ill{} \ill{} \ill{}
\uncertain{est} \uncertain{loc.}\ \uncertain{triviale}.
\note{deux lignes interlinéaires serrées, en partie biffées, et une
addition au-dessus de la ligne (« \uncertain{à} \uncertain{priori} ») ; la
fin de la phrase est très incertaine}

Si $X = \operatorname{Spec}(A)$ \uncertain{est} \uncertain{le}
\uncertain{spectre} d'un \ill{} \emph{\uncertain{local}}, \uncertain{on}
\uncertain{aura} :
\[
\begin{array}{ll}
H^{i}(X, G) = 0 & \text{si } i > 0 \\
H^{i}(X - a, G) = 0 & \text{si } i \geq \dim A
\end{array}
\]
\note{après « $\dim A$ », un groupe biffé où se lit $(X - a) + 1$ ; au-dessous,
entre crochets : « [\uncertain{car} \uncertain{dim.}\ \uncertain{comb.}\
$(X - a) \leq \dim A - 1$] »}
d'où
\[
0 \to H^{0}(X/X - a, G) \to H^{0}(X, G) \to H^{0}(X - a, G) \to H^{1}(X/X - a, G) \to 0
\]
\[
H^{i}(X/X - a, G) \overset{\partial}{\longleftarrow} H^{i-1}(X - a, G) \qquad (i \geq 2)
\]
[\emph{\uncertain{c'est} \uncertain{nul} si $i > \dim A$}]
\note{la flèche $\partial$ est double, $\leftleftarrows$ ou $\simeq$ ; le
dernier signe entre crochets est surchargé}

\uncertain{On} \uncertain{suppose} $A$ \emph{\uncertain{régulier}}.

a)
\[
H^{0}_{A}(G) =
\left\lbrace
\begin{array}{ll}
G(A) = G(K) & \text{si } \dim A = 0 \\
0 & \text{si } \dim A > 0
\end{array}
\right.
\]
Le \uncertain{foncteur} $H^{0}_{A}(G)$ \uncertain{est}
\uncertain{exact} \uncertain{pour} les \uncertain{suites}
\struck{\ill{}} \uncertain{loc.}\ \uncertain{triviales}
\struck{(\ill{} \ill{} \ill{} $A$ \ill{})}.
\note{au-dessus du mot biffé, un mot ajouté, peut-être
« \uncertain{courtes} »}

b)
\[
H^{1}_{A}(G) =
\left\lbrace
\begin{array}{ll}
0 & \text{si } \dim A \neq 1 \\
G(K)/G(A) & \text{si } \dim A = 1
\end{array}
\right.
\]
\note{à droite de la première ligne, une parenthèse biffée et surchargée,
où se lisent « \uncertain{distinguer} » et $A$}
Le \uncertain{foncteur} $H^{1}_{A}(G)$ \uncertain{est}
\uncertain{exact} \uncertain{pour} les \uncertain{suites}
\uncertain{loc.}\ \uncertain{triviales}.
\marginal{\uncertain{Prouver} \uncertain{que} si $\dim A = 1$,
$J_{A/k}$ \uncertain{est} \emph{\uncertain{connexe}}, \uncertain{sans}
\uncertain{composante} \uncertain{abélienne} (\uncertain{car}
$H^{1}_{A}(G)$ \uncertain{est} \uncertain{nul} \uncertain{pour} $G$
\uncertain{propre} \uncertain{ou} \uncertain{discret})}
\note{la note marginale est écrite en oblique dans la marge gauche ; le
$J$ de $J_{A/k}$ est surchargé}

c)
\[
H^{2}_{A}(G) = H^{1}(X - a, G) =
\left\lbrace
\begin{array}{l}
0 \quad \text{si } \dim A \neq 2 \\
\text{\ill{}}
\end{array}
\right.
\]
\ill{} \ill{} \ill{} \ill{} \uncertain{groupe} \ill{} \ill{}
\uncertain{composantes} \uncertain{connexes} \ill{} \uncertain{de} $G$.
\note{dans $H^{1}(X - a, G)$, l'exposant $1$ est écrit sur un $2$ ; à
droite de la première ligne, un passage lourdement biffé, où se lit
« $\mathrm{Ab}$ »}

\page{71}
\note{la page est presque blanche ; au bas, tête-bêche par rapport au
numéro de l'archiviste, neuf lignes qui donnent une autre rédaction des
points a)--c) de la page 70 (probablement une première version). Un long
trait vertical, peut-être une annulation, traverse les formules}

a)
\[
H^{0}_{A}(G) = \mathrm{Ker}\bigl[G(\operatorname{Spec}(A)) \to G(\operatorname{Spec}(A) - a)\bigr]
=
\left\lbrace
\begin{array}{ll}
G(A) & \text{si } \dim A = 0 \\
0 & \text{si } \dim A \neq 0
\end{array}
\right.
\]
\uncertain{Si} $\dim A \neq 0$, \uncertain{le} \uncertain{foncteur}
$H^{0}_{A}(G)$ \uncertain{est} \uncertain{exact} \uncertain{pour} les
\uncertain{suites} \uncertain{exactes} \uncertain{qui} \uncertain{sont}
\uncertain{loc.}\ \uncertain{triviales}.
\note{devant cette ligne, dans la marge, un « b » \uncertain{surchargé}}

b)
\[
H^{1}_{A}(G) = \mathrm{Coker}\bigl(G(\operatorname{Spec}(A)) \to G(\operatorname{Spec}(A) - a)\bigr)
=
\left\lbrace
\begin{array}{ll}
G(K)/G(A) & \text{si } \dim A = 1 \\
0 & \text{si } \dim A \neq 1
\end{array}
\right.
\]

c)
\[
H^{2}_{A}(G) = H^{1}_{?}(\operatorname{Spec}(A) - a, G)
\]
\uncertain{il} \uncertain{faut} \uncertain{choisir} \uncertain{sa}
\uncertain{théorie} : \ill{} \ill{} (\uncertain{fibrés}
\uncertain{loc.}\ \uncertain{triviaux}, \uncertain{strictement}
\uncertain{isotriviaux}, \uncertain{isotriviaux} \ldots). \uncertain{On}
\uncertain{trouve} :
\note{l'exposant de $H^{2}_{A}$ est d'une lecture incertaine ; l'indice
de $H^{1}$ est un « ? » ; la page s'arrête sur « \uncertain{On}
\uncertain{trouve} : »}

\page{72}
\note{suite de la page 70 (point c)). Il écrit deux sortes de $H$ : le
$H$ droit de la théorie de Zariski, et un $H$ cursif, rendu ici
$\mathcal{H}$, pour la théorie des fibrés principaux quelconques}
\uncertain{Si} $\dim A = 0$ ou $1$, il \uncertain{n'y} \uncertain{a}
\uncertain{rien} \uncertain{à} \uncertain{prouver}, \uncertain{on}
\uncertain{peut} \uncertain{supposer} $\dim A \geq 2$.
\uncertain{Soit} $\mathcal{H}^{1}$ \uncertain{la} \struck{\ill{}}
\uncertain{théorie} des \uncertain{classes} \uncertain{de}
\uncertain{fibrés} \uncertain{principaux} \uncertain{homogènes}
(\uncertain{triviaux} \uncertain{par} \uncertain{extension}
\uncertain{fid.}\ \uncertain{plate} \uncertain{quasi-compacte}
\uncertain{de} \uncertain{la} \uncertain{base}). \uncertain{On}
\uncertain{aura} $H^{1}(X - a, G) \subset \mathcal{H}^{1}(X - a, G)$,
\uncertain{qui} \uncertain{lui} \uncertain{est} \uncertain{égal} \ill{}
\ill{} \uncertain{produit} \ill{} \ill{} $G$.
\uncertain{On} \uncertain{a} si \struck{$\dim A =$} $G$ \uncertain{est}
\ill{} \ill{}
\struck{\ill{} \ill{} \ill{}}
\note{au-dessus du passage biffé, « \uncertain{propre} \uncertain{sur}
\uncertain{la} » ; en tête de page, à droite, « $\dim A = 2$ » d'une autre
encre}
\[
H^{1}(X - a, G) = 0 \qquad (\text{\uncertain{loc.}\ \uncertain{trivial}} = \text{\uncertain{trivial}}),
\]
\note{avant « $= 0$ », un groupe lourdement biffé}
\uncertain{donc} \ill{} \ill{} \uncertain{tout} \ill{} \ill{}
$H^{1}(X - a, G)$ \uncertain{provient} \uncertain{de}
$\mathcal{H}^{1}(X - a, L)$, $L$ \uncertain{le} \uncertain{sous-groupe}
\uncertain{linéaire} [\uncertain{connexe}]
\note{« [connexe] » est ajouté au-dessus de « linéaire », entre crochets}
\uncertain{On} \uncertain{peut} \uncertain{même} \uncertain{prendre}
\uncertain{linéaire} \uncertain{connexe}, \ill{} \ill{} \uncertain{sans}
\uncertain{que} \struck{$H^1 = \mathcal{H}^1$ (}
$H^{1}(G) = \mathcal{H}^{1}(G)$, \uncertain{donc} \ill{} \ill{}
\uncertain{aura}
\[
H^{1}(X - a, G) = \mathrm{Im}\bigl(H^{1}(X - a, L) \to H^{1}(X - a, G)\bigr)
\]
\[
\mathrm{Im}\bigl(\mathcal{H}^{1}(X - a, L) \to \mathcal{H}^{1}(X - a, G)\bigr) = H^{1}(X - a, L) / \mathrm{Im}\, H^{0}(X - a, G/L)
\]
\note{dans la première ligne, le terme biffé est
$H^{1}(X - a, L)$ ; le « $\mathrm{Im}$ » du début de la seconde est biffé
d'un trait}
\uncertain{Or} $H^{0}(X - a, G/L) = H^{0}(X, G/L)$ \uncertain{puisque}
\ill{} \ill{} \ill{},
$H^{1}(X, L) = 0$, \uncertain{donc}
$\underline{H^{1}(X - a, G) = H^{1}(X - a, L)}$.
\uncertain{D'ailleurs}, \uncertain{comme} $X$ \uncertain{est}
\uncertain{factoriel}, $H^{1}(X - a, \mathbb{G}_m) = 0$,
\uncertain{d'où} $H^{1}(X - a, L) = H^{1}(X - a, U)$, $U$
\uncertain{le} \uncertain{composant} \uncertain{unipotent}.
\uncertain{Cette} \uncertain{dernière} \uncertain{se} \uncertain{ramène}
\uncertain{à} \uncertain{les} $H^{1}(X - a, \mathbb{G}_a)$,
\uncertain{qui} \uncertain{sont} \uncertain{nuls} \uncertain{si}
$\dim A > 2$, \uncertain{mais} \uncertain{non} \uncertain{pour}
\ill{} $\dim A = 2$.
\note{« $\mathbb{G}_m$ » et « $\mathbb{G}_a$ » sont lus sur des $G$ à
indice peu net}

d) $H^{3}_{A}(G) = H^{2}(X - a, G)$
\uncertain{On} \uncertain{fait} \uncertain{généralement},
\ill{} \struck{$H^{i}(X - a, G) =$} \ill{} \ill{} \ill{} \ill{}
\uncertain{de} \uncertain{faisceaux}
\note{l'exposant de $H_A$ est d'une lecture incertaine ; on attend 3}
\[
0 \to \underline{O}_{x}(G) \to \underline{R}_{x}(G) \to \underline{D}_{x}(G) \to 0
\]
d'où, \uncertain{puisque} $\underline{R}_{x}(G)$ \uncertain{est}
\uncertain{flasque},
\[
H^{i}\bigl(\underline{O}_{x}(G)\bigr) = H^{i-1}\bigl(\underline{D}_{x}(G)\bigr) \quad \text{si } i \geq 2
\]
\note{la dernière ligne est écrite d'une autre encre, serrée au pied de la
page}

\page{73}
D'\uncertain{autre} \uncertain{part}, \uncertain{comme} si $L$
\uncertain{est} \uncertain{la} \uncertain{partie} \uncertain{linéaire}
\uncertain{de} $G$, \uncertain{on} \uncertain{a}
$\underline{D}_{x}(G) \simeq \underline{D}_{x}(L)$, d'où
\note{avant $\underline{D}_x(G)$, un $H$ biffé}
\[
H^{i}(X - a, G) \simeq H^{i}(X - a, L)
\qquad
\Bigl(i \geq 2,\ \text{\uncertain{et}}\ = \text{\ill{} \ill{} \uncertain{pour}}\ i \geq 1\Bigr)
\]
\uncertain{On} \uncertain{sait} \uncertain{que} \uncertain{l'on} \ill{}
$H^{i}(X - a, \mathbb{G}_m) = 0$ \uncertain{puisque} $X$ \uncertain{est}
\uncertain{factoriel}, \uncertain{donc} \uncertain{on} \uncertain{aura}
\[
H^{i}(X - a, G) = H^{i}(X - a, U) \qquad (i \geq 1)
\]
\struck{\ill{}} D'\uncertain{ailleurs}, \uncertain{on} \uncertain{a}
\[
H^{i}(X - a, \mathbb{G}_a) = 0 \quad \text{si } i \neq \dim A - 1
\]
\uncertain{donc}
\[
H^{i}(X - a, U) = 0 \quad \text{si } i \neq \dim A - 1
\]
\uncertain{Cela} \uncertain{achève} \uncertain{de} \uncertain{prouver}
\ill{} \ill{} \ill{}.
\note{dans $H^{i}(X - a, \mathbb{G}_a)$, le $\mathbb{G}_a$ est écrit sur
une autre lettre}

\uncertain{Noter} \uncertain{aussi}
$H^{i}_{A}(G) \xrightarrow{\ \sim\ } H^{i}_{\hat{A}}(G)$ :
\uncertain{rien} \uncertain{ne} \uncertain{change} \uncertain{en}
\uncertain{remplaçant} $A$ \uncertain{par} $\hat{A}$.

\page{74}
\emph{\uncertain{Jacobienne}} \uncertain{locale} [\uncertain{relative}]
\uncertain{d'une} \uncertain{localité} (\uncertain{régulière}) $A$
\uncertain{sur} \uncertain{un} \uncertain{corps} $k$.
\note{« relative » et « régulière » sont ajoutés au-dessus de la ligne}
$J_{A/k}$ \uncertain{est} \uncertain{défini} \uncertain{par} \uncertain{le}
\uncertain{foncteur} $\mathrm{Hom}_{\text{\uncertain{k-gr.}}}(J_{A/k}, G)$
\uncertain{en} $G$ (\uncertain{groupe} [\uncertain{loc.}] \uncertain{de}
\uncertain{type} \uncertain{fini} \uncertain{sur} $k$) \uncertain{on}
\uncertain{doit} \uncertain{avoir}
\[
\mathrm{Hom}_{\text{k-gr}}(J_{A/k}, G) \simeq H^{i}_{A}(G) \qquad i = \dim A
\]
$H^{i}_{A}$ \uncertain{est} \uncertain{le} \uncertain{foncteur}
\uncertain{cohomologique} \uncertain{de} \uncertain{localisation}
\uncertain{par} \uncertain{l'anneau} \uncertain{local} $A$
[\uncertain{de} $\operatorname{Spec}(A)$ \uncertain{à}
$\operatorname{Spec}(A) - \mathfrak{m}_{A}$] \uncertain{foncteur}
\uncertain{qui} \uncertain{est} \uncertain{défini} \uncertain{à}
\uncertain{priori} \uncertain{sur} \uncertain{les} \uncertain{faisceaux}
\uncertain{de} Weil \uncertain{sur} $\operatorname{Spec}(A)$.
\uncertain{On} \uncertain{a}
\[
\begin{array}{lll}
H^{0}_{A}(G) = G(K) & & \text{si } A = K \text{ de dim.\ } 0 \\
H^{1}_{A}(G) = G(K)/G(A) & & \text{si } A \text{ de dim.\ } 1,\ \text{de corps des fractions } K \\
H^{i}_{A}(G) = H^{i-1}(\operatorname{Spec}(A) - \mathfrak{m}_{A}, G) & & \text{si } A \text{ de dim.\ } \geq 2
\end{array}
\]
\marginal{\uncertain{Le} \uncertain{foncteur} \uncertain{en} $G$
\uncertain{est} \uncertain{représentable}, \uncertain{vrai} \ill{}
\uncertain{fait} \uncertain{pour} $H^{i}_{A}(G) = 0$ \uncertain{si}
$A$ \uncertain{est} \emph{\uncertain{régulier}}, \uncertain{bon}
\uncertain{pour} \ill{}}
\marginal{\uncertain{Tout} \uncertain{le} \ill{} \uncertain{est}
\uncertain{des} $J_{A/k}$ \uncertain{proalgébriques}
\uncertain{commutatifs} \uncertain{sur} $k$.}
\note{deux notes marginales en oblique dans la marge gauche, en haut}

\uncertain{Je} \uncertain{soupçonne} :

a) $J_{K/k}$ \uncertain{est} \uncertain{une} \emph{\uncertain{extension}
\uncertain{de} $\mathbb{Z}_{(k)}$} (\uncertain{qui} \uncertain{correspond}
\uncertain{au} \uncertain{foncteur} $G \mapsto G(k)$) \uncertain{par}
\uncertain{un} \uncertain{groupe} \uncertain{connexe}, \uncertain{qui}
\uncertain{a} \uncertain{lui-même} \emph{\uncertain{une} \uncertain{partie}
\uncertain{abélienne}} (\uncertain{c'est} \uncertain{la}
\uncertain{Albanese} $M$ \uncertain{de} $K/k$), \emph{\uncertain{une}
\uncertain{partie} \uncertain{multiplicative}}, \uncertain{et}
\uncertain{la} \uncertain{partie} \uncertain{unipotente} $U$ \ill{}
$J_{K/k}$ \ill{} \uncertain{mais} \uncertain{son} $\mathrm{Ext}^{1}$
\ill{} \struck{\ill{} \ill{}} \uncertain{non} \uncertain{nul}
\uncertain{en} \uncertain{général} \ill{} \ill{} \ill{}.
\note{la fin de a) est serrée en interligne, avec des mots biffés et
d'autres ajoutés ; « $M$ » et « $U$ » sont écrits au-dessus des mots
qu'ils nomment. Une ligne courbe relie « $\mathbb{Z}_{(k)}$ » à une note
serrée, verticale, dans la marge droite, illisible}

b) \uncertain{Si} $A$ \uncertain{est} \uncertain{de} \uncertain{dim.}\ 1,
$J_{A/k}$ \uncertain{est} \emph{\uncertain{connexe}},
\emph{\uncertain{linéaire}} (\uncertain{n'a} \uncertain{pas}
\uncertain{de} \uncertain{partie} \uncertain{abélienne}) \uncertain{mais}
\uncertain{elle} \struck{\ill{} \ill{}} \uncertain{a} \struck{\ill{}}
\uncertain{une} \uncertain{partie} \uncertain{multiplicative}
$M \simeq \mathbb{G}_m$.
\marginal{N.B. $D(M) = \mathrm{Hom}_{\text{k-gr}}(M, \mathbb{G}_m)$
= \uncertain{groupe} \uncertain{des} \uncertain{diviseurs}
\struck{\ill{}} \ill{} $\geq 0$ \ill{} (\uncertain{ou} \ill{}
\uncertain{diviseurs} \ill{} $K/k$) \ill{}
$\mathrm{Hom}(U, \mathbb{G}_a)$ \ill{} \ill{} \ill{} \ill{} \ill{}}
\note{la note marginale du bas de la marge gauche est écrite en oblique,
très surchargée et en partie biffée ; seuls quelques termes se lisent}

\page{75}
\uncertain{De} \uncertain{plus}, $J_{A/k}$ \uncertain{est}
\emph{\uncertain{projectif}} \uncertain{relativement} \uncertain{aux}
\uncertain{suites} \uncertain{exactes} \emph{\uncertain{loc.}\
\uncertain{triviales}}, \uncertain{i.e.}\ \uncertain{aux}
\emph{\uncertain{extensions}} \emph{\uncertain{par} \uncertain{un}
\uncertain{groupe} \uncertain{linéaire} \uncertain{connexe}}
\ill{} $0$.
\note{ces deux lignes, serrées, sont ajoutées en tête de page, au-dessus
de c)}

c) \uncertain{Si} $\dim A \geq 2$, \uncertain{alors} $J_{A/k}$
\uncertain{est} \emph{\uncertain{connexe}}, \emph{\uncertain{unipotent}}
(\uncertain{i.e.}\ \uncertain{plus} \uncertain{de} \uncertain{partie}
\uncertain{abélienne} \uncertain{ni} \uncertain{multiplicative}) et
\uncertain{de} \uncertain{plus} \emph{\uncertain{projectif}}
(\uncertain{car} \uncertain{le} \uncertain{foncteur} \uncertain{qu'il}
\uncertain{représente} \uncertain{est} \uncertain{exact}),
\struck{\uncertain{comme} \uncertain{pour} $\dim A = 1$}.
\note{sous le texte, un paraphe ; le reste de la page est blanc}

\page{76}
\note{la section II fait suite au « I) Théorie Čechiste » de la
page 70}

II \uncertain{Cohomologie} \uncertain{de} Weil \ldots\
[\uncertain{semble} \uncertain{marcher} \uncertain{moins}
\uncertain{bien} \ldots]

\uncertain{Je} \uncertain{crois} \uncertain{que} \uncertain{vrai}
\uncertain{en} \uncertain{tout} \uncertain{cas} [\uncertain{pour}
$A$ \uncertain{local} \uncertain{hensélien}] \uncertain{que}
\[
H^{i}(A, G) = 0 \quad \text{si } i > 0.
\]
\struck{\ill{}} \uncertain{Si} $G$ \uncertain{est} \uncertain{$k$-simple},
\uncertain{alors} $H^{i}(A, G) = H^{i}(L, G)$ [\uncertain{si} $i > 0$],
$L$ \uncertain{étant} \uncertain{le} \uncertain{corps} \uncertain{résiduel},
\uncertain{et} \struck{\ill{} \ill{}} \ill{} \ill{}
\[
H^{i}(L, G) = H^{i}(L, G/L_{0})
\]
\uncertain{où} $L_{0}$ \uncertain{est} \uncertain{le} \uncertain{plus}
\uncertain{grand} \uncertain{sous-groupe} \uncertain{linéaire}
\uncertain{connexe} \uncertain{de} $G$.
\note{« [pour $A$ local hensélien] » et « si $i > 0$ » sont ajoutés
au-dessus des lignes et rattachés par des traits ; la lecture de
$L_{0}$ (un $\mathcal{L}$ bouclé) est incertaine}

\struck{\ill{} \ill{} \ill{} \ill{} \ill{}}
\struck{$H^{1}(X/X - a, G) \to H^{1}(X, G)$ \ill{} \ill{}
\uncertain{isomorphisme}} \struck{\uncertain{si} $\dim A = 0$,
\ill{} \ill{} \ill{}} \struck{\uncertain{si} $\dim A = 1$, \ill{}}
$H^{1}(X, G)$, \ill{} \ill{} \ill{}, \ill{} \ill{} \ill{}
\note{tout ce bloc, jusqu'à la formule suivante, est encadré et barré
d'une grande ligne ondulée ; sous $H^{1}(X, G)$, « \uncertain{e.h.} »}
\[
H^{1}_{A}(G) = \mathrm{Coker}\bigl(H^{0}(X) \to H^{0}(X - a)\bigr)
\]

\uncertain{Si} $\dim A = 0$, \uncertain{rien} \uncertain{à}
\uncertain{prouver}.

\uncertain{Si} $\dim A \geq 1$, \uncertain{alors}
$H^{1}(X/X - a, G)$ \uncertain{est} \uncertain{extension}
\uncertain{de}
\[
\mathrm{Ker}\bigl(H^{1}(X, G) \to H^{1}(X - a, G)\bigr)
\quad \text{\uncertain{par}} \quad G(K)/G(A).
\]
\uncertain{Or}
\[
\mathrm{Ker}\bigl(H^{1}(X, G) \to H^{1}(X - a, G)\bigr) = 0
\]
\uncertain{si} $G$ \uncertain{est} \uncertain{simple}, \uncertain{car}
\struck{\ill{} \ill{} \ill{} \ill{}} \uncertain{car} \uncertain{si}
$G$ \uncertain{est} \uncertain{extension} \uncertain{d'un}
\uncertain{groupe} \uncertain{discret} \uncertain{par} \uncertain{un}
\uncertain{groupe} \uncertain{propre}, \uncertain{mais} \uncertain{dans}
\uncertain{un} \uncertain{cas} \ill{} \uncertain{les} \ill{}
\uncertain{d'un} \uncertain{fibré} \uncertain{principal}
\uncertain{est} \uncertain{partout} \uncertain{défini}
[\uncertain{en} \uncertain{codim.}\ \uncertain{au} \uncertain{moins} 1,
\uncertain{donc} \uncertain{partout}].
\note{en marge gauche, un petit diagramme :}
\[
\begin{array}{ccc}
0 & & \\
\downarrow & & \\
H^{1}(X, G) & \longrightarrow & H^{1}(X - a, G) \\
\downarrow & & \downarrow \\
H^{1}(X, G/L) & \longrightarrow & H^{1}(X - a, G/L)
\end{array}
\qquad \text{\uncertain{si}}
\]
\note{« \uncertain{Or} » est écrit au-dessus d'un mot biffé, en face de
ce diagramme}

\page{77}
\note{le signe entre $A$ et $K$ dans $H^{i}(A \bmod K, G)$ est un mot
bref, lu « mod » sans certitude : la cohomologie relative de $A$ modulo
son corps des fractions}
\uncertain{Utilisant} \ill{} \ill{}
$H^{1}_{A}(G) = G(K)/G(A)$, \uncertain{vrai} \ill{} \ill{} \ill{}
$\dim A \geq 1$, \ill{}

[\uncertain{Supposons} \emph{$\dim A = 1$}. \uncertain{On}
\ill{} \ill{} $H^{i}_{A}(G) = 0$ \ldots
\note{l'indice de ce $H$ est d'une lecture incertaine}
\[
\to H^{1}(A, G) \to H^{1}(K, G) \xrightarrow{\ \partial\ } H^{2}(A \bmod K, G) \to H^{2}(A, G)
\]
\[
\to H^{2}(K, G) \xrightarrow{\ \partial\ } H^{3}(A \bmod K, G)
\]
\note{la dernière flèche descend verticalement, au bout de la ligne}
\uncertain{donc} \uncertain{on} \uncertain{doit} \uncertain{trouver}
\uncertain{les} $H^{i}(A \bmod K, G) = H^{i}_{A}(G)$ \uncertain{par}
\uncertain{extension} \uncertain{d'un} \uncertain{noyau} \uncertain{par}
\uncertain{un} \uncertain{conoyau} \uncertain{de}
\[
\begin{array}{ccc}
H^{i}(A, G) & \longrightarrow & H^{i}(K, G) \\
\| & & \\
H^{i}(L, G) & &
\end{array}
\]

\struck{Si $G$ est}
\uncertain{Noter} \uncertain{que} \uncertain{si} $L_{0}$ \uncertain{est}
\uncertain{le} \uncertain{plus} \uncertain{grand}
\uncertain{sous-groupe} \uncertain{linéaire} \uncertain{connexe} \uncertain{de}
$G$, \uncertain{alors}
\[
\left.
\begin{array}{l}
H^{i}(A, G) = H^{i}(A, G/L_{0}) \\
H^{i}(K, G) = H^{i}(K, G/L_{0})
\end{array}
\right| \ i \geq 1
\]
\uncertain{donc} \uncertain{on} \uncertain{peut} \uncertain{supposer}
$L_{0} = 0$.]


\[
\begin{array}{ccc}
\pi_{1}(L) = \pi_{1}(A) & \pi' & B_{\pi'} \\
\uparrow & & \uparrow \\
\pi_{1}(K) & \pi & B_{\pi} \\
\uparrow & & \uparrow \\
\pi_{1}(A \bmod K) & \pi'' & B_{\pi''} \\
\uparrow & & \uparrow \\
e & & \ill{}
\end{array}
\]
\note{au-dessus de $\pi_{1}(L)$, une flèche montante vers « $e$ » ; le
tableau s'arrête au bas de la page}

\page{78}
\emph{\uncertain{Supposons} $\dim A \geq 2$}, \uncertain{on}
\uncertain{a} \uncertain{déjà} $H^{i}_{A}(G) = 0$ \uncertain{si}
$i = 0, 1$.

\struck{Alors} $H^{2}(X/X - a, G)$ \uncertain{est} \uncertain{une}
\uncertain{extension} \uncertain{de}
$\mathrm{Coker}\bigl(H^{1}(X, G) \to H^{1}(X - a, G)\bigr)$
\uncertain{par} $\mathrm{Ker}\bigl(H^{2}(X) \to H^{2}(X - a, G)\bigr)$.
\uncertain{Si} $G$ \uncertain{est} \uncertain{extension} \uncertain{d'un}
\uncertain{discret} \uncertain{par} \uncertain{un} \uncertain{propre},
[\uncertain{et} \uncertain{est} \emph{\uncertain{simple}} \uncertain{sur}
$k$], \uncertain{alors} \uncertain{par} \uncertain{le}
\emph{\uncertain{théorème} \uncertain{de} \uncertain{pureté}}
$H^{1}(X, G) \to H^{1}(X - a, G)$ \uncertain{est}
\emph{\uncertain{surjectif}} [\uncertain{et} \uncertain{si} $G$
\uncertain{est} \ill{} ? \ldots]. \uncertain{Quand} \ill{}
$H^{2}(X) \to H^{2}(X - a, G)$
\note{« [et est simple sur $k$] » est ajouté au-dessus de la ligne ; la
page s'arrête sur cette formule, le reste est blanc}

\page{79}
\marginal{\uncertain{Adèles} \uncertain{et} \uncertain{Jacobiennes}
\uncertain{généralisées} \uncertain{des} \uncertain{courbes}}
\note{la note marginale, en haut à gauche, est écrite en oblique et
encadrée : elle sert de titre à la page}
\[
\underline{\underline{\mathrm{Ext}}}^{*}(J_{*}, G)
\Longleftarrow H^{p}\bigl(\mathrm{Ext}^{q}(J_{*}, G)\bigr)
\]
\[
\qquad\qquad\qquad\ \Longleftarrow \mathrm{Ext}^{p}\bigl(H_{q}(J_{*}), G\bigr)
\]
\note{sous $\underline{\underline{\mathrm{Ext}}}^{*}(J_{*}, G)$, « $\overset{?}{=}$ »
et un $H^{*}(X, G)$ barré d'une croix, avec, entouré, « \uncertain{à}
\uncertain{voir} ». Les termes $\mathrm{Ext}^{q}(J_{*}, G)$ et
$H_{q}(J_{*})$ sont entourés de deux ovales sécants, le premier marqué
« ? » ; une flèche mène à une note de la marge droite :}
\marginal{\uncertain{cette} \uncertain{suite} \uncertain{spectrale}
\uncertain{exprime} \uncertain{un} \uncertain{lien} \ill{} \ill{}
\uncertain{de} Weil \uncertain{et} \uncertain{coh.}\ Č\uncertain{echiste}
\ill{}, \ill{} \uncertain{localisant} \ldots}
\note{sous les deux ovales, une ligne biffée ; au-dessus, ajouté :
« \uncertain{i.e.}\ \ill{} \ill{} \uncertain{inévitable}
\ill{} \ill{} \ill{} ?? \uncertain{invrai} »}

$X$ \uncertain{de} \uncertain{dim.}\ \textbf{1} (\uncertain{pas}
\uncertain{nécessairement} \uncertain{complète}, \uncertain{mais}
\uncertain{sans} \uncertain{singularités})
\[
H_{0}(J_{*}) = J^{\bullet}_{X/k}
\]
\uncertain{la} \uncertain{jacobienne} \uncertain{de} Rosenlicht
\uncertain{de} $X$, \uncertain{extension} \uncertain{de} $\mathbb{Z}$
\uncertain{par} \uncertain{l'Albanese} \uncertain{birationnel},
\uncertain{par} \uncertain{le} \uncertain{produit} \uncertain{des}
\uncertain{groupes} \uncertain{multiplicatifs} \uncertain{des}
\uncertain{points} \ill{} \uncertain{qui} \uncertain{manquent} \uncertain{à}
$X$ \uncertain{pour} \uncertain{être} \uncertain{complète}, \ill{}
\uncertain{divisé} \uncertain{par} $\mathbb{G}_m$ \uncertain{diagonal}
\note{« \uncertain{des} \uncertain{points} » est ajouté au-dessus de la ligne}
\[
H_{1}(J_{*}) =
\left\lbrace
\begin{array}{ll}
\mathbb{G}_m & \text{si } X \text{ complète} \\
0 & \text{si } X \text{ non complète}
\end{array}
\right.
\]
\uncertain{Donc} \uncertain{si} $X$ \uncertain{est} \uncertain{non}
\uncertain{complète}, \uncertain{on} \uncertain{a}
\[
H^{*}(X, G) = \underline{\underline{\mathrm{Ext}}}^{*}(J_{*}, G) = \mathrm{Ext}^{*}(J^{\bullet}_{X/k}, G)
\]
\note{le premier terme est écrit sur un autre, lu $H^{*}(X, G)$}
\uncertain{Si} $X$ \uncertain{est} \uncertain{complète}, \uncertain{on}
\uncertain{a} \uncertain{une} \uncertain{suite} \uncertain{exacte}
\[
\cdots \to \mathrm{Ext}^{p}(J^{\bullet}_{X/k}, G) \to H^{p}(X, G) \to \mathrm{Ext}^{p-1}(\mathbb{G}_m, G) \to \mathrm{Ext}^{p+1}(J^{\bullet}_{X/k}, G)
\]
\note{la dernière flèche porte, sous elle : « \uncertain{multiplication}
\uncertain{par} \uncertain{une} \uncertain{classe} \uncertain{de}
$\mathrm{Ext}^{2}(J^{\bullet}_{X/k}, \mathbb{G}_m)$ » ; dans $H^{p}(X, G)$,
le $X$ est surchargé}

\uncertain{Passant} : \uncertain{la} \uncertain{limite}
\uncertain{inductive} \uncertain{sur} \uncertain{les} \uncertain{ouverts}
\uncertain{denses} \uncertain{de} $X$, \uncertain{on} \uncertain{trouve}
\uncertain{une}, \uncertain{si} $K = k(X)$
\[
\boxed{\mathrm{Ext}^{*}(J^{\bullet}_{K/k}, G) \simeq H^{*}(K, G)}
\]

\page{80}
\uncertain{Calcul} \uncertain{de} $J^{\bullet}_{A/k}$ \uncertain{si}
$A$ \uncertain{est} \uncertain{un} \uncertain{anneau} \uncertain{de}
\uncertain{val.}\ \uncertain{discrète} \uncertain{dont} \uncertain{le}
\uncertain{corps} \uncertain{résiduel} $L$ \uncertain{est} \uncertain{de}
\uncertain{degré} \uncertain{de} \uncertain{transcendance} 1
\note{en tête, un signe lourdement noirci ; dans $J^{\bullet}_{A/k}$, le
$A$ est écrit sur une autre lettre. $\coprod$ rend son produit renversé,
$\mathbb{S}$ son $S$ barré (comme au lot 2)}
\[
\coprod_{L/k} \mathbb{S}^{*}_{L}
= \underbrace{\coprod_{L/k} \mathbb{Z}_{L}}_{J^{*}_{L/k}}
\times \underbrace{\coprod_{L/k} \mathbb{G}_{m\,L}}_{\mathbb{G}_{m\,k}}
\times \coprod_{L/k} \mathbb{S}'_{L}
\]
\[
\mathcal{U} = \coprod_{L/k} \mathbb{S}'_{L}
\quad \text{\uncertain{est} \uncertain{un} \uncertain{groupe}
\emph{\uncertain{unipotent}} \emph{\uncertain{connexe}}.}
\]
\[
\mathrm{Hom}_{\text{k-gr}}(\mathcal{U}, G)
= \varinjlim_{n} \mathrm{Hom}_{\text{L-gr}}(\mathbb{S}'^{(n)}_{L}, G_{L})
= \varinjlim_{n} \varinjlim_{C'} \mathrm{Hom}_{\text{C'-gr}}(\mathbb{S}'^{(n)}_{C'}, G_{C'})
\]
\note{sous $\varinjlim_{C'}$ : « \ill{} \ill{} \uncertain{ouverts}
\uncertain{de} $C$ »}
\[
= \varinjlim_{C'} \varinjlim_{n} \mathrm{Hom}_{\text{C'-gr}}(\mathbb{S}'^{(n)}_{C'}, G_{C'})
\]
\note{sous cette ligne, une accolade vers le terme suivant :}
\[
\mathrm{Hom}_{\substack{\text{C'-schémas} \\ \text{\uncertain{à} \uncertain{sections}} \\ \text{\uncertain{marquées}}}}\bigl(E^{1}_{C'}, G_{C'}\bigr)
\quad \text{\uncertain{puis} \uncertain{à} Rosenlicht}
\]
\[
\|
\]
\[
\mathrm{Hom}^{\text{\uncertain{pointés}}}_{E^{1}\text{-schémas}}\bigl(C'_{E^{1}}, G_{E^{1}}\bigr)
\]
\marginal{\uncertain{envoyant} \uncertain{ceux} \uncertain{qui}
\struck{\ill{}} \ill{} \uncertain{la} \uncertain{section} $0$
\uncertain{sur} $C'_{E^{1}} \times k(\alpha)$ [\uncertain{le}
\uncertain{point} \uncertain{à} \uncertain{l'infini} \uncertain{de}
$E^{1}$]}
\note{une flèche courbe relie le second $\mathrm{Hom}$ à la note de la
marge droite}
\[
\|
\]
\[
\mathrm{Hom}^{\text{\uncertain{pointés}}}_{E^{1}\text{-gr}}\bigl(\underbrace{J^{\bullet}_{C'_{E^1}/E^1}}_{J^{\bullet}_{C'/k} \times E^{1}}, G_{E^{1}}\bigr)
\]
\[
\|
\]
\[
\mathrm{Hom}^{\text{\uncertain{restreint}}}\bigl(J^{\bullet}_{C'/k} \times E^{1}, G\bigr)
\]
\note{sous « restreint », un mot biffé, puis
« \uncertain{multiplicatif} \uncertain{en} \uncertain{les}
\uncertain{facteurs} \uncertain{k-gr.} » ; à gauche, un petit schéma :
$E^{1}_{C'} = E^{1} \times C' = C'_{E^{1}}$, une flèche vers $E^{1}$,
$G_{E^{1}} = E^{1} \times G$, et $E^{1} \times C' \to G \times C'$}
\[
\mathrm{Hom}_{\text{k-gr}}(\mathcal{U}, G)
= \mathrm{Hom}^{\text{\uncertain{restreint}}}\bigl(J^{\bullet}_{L/k} \times E^{1}, G\bigr)
\]
\note{sous « restreint » : « \uncertain{multiplicatif} \uncertain{par}
\uncertain{rapport} \ill{} \ill{} »}
\[
0 \to \mathcal{U}_{0} \to \mathcal{U} \to \mathbb{S}'_{k} \to 0,
\quad \text{\uncertain{donc}}
\]
(\uncertain{suite} \uncertain{splittante})
\note{la page s'arrête sur « donc »}
\note{la page 81 (lot suivant) semble continuer cette ligne :
$\mathrm{Hom}_{k\text{-gr}}(\mathcal{U}_0, G) = \mathrm{Hom}(\ldots)$}

\end{document}
